\chapter{Hamiltonian Semantic Fields}

\section{What Gradient Descent Cannot Explain}

Chapter 19 left an honest debt. Gradient descent on H explains why a field settles, why coherence smooths, why entropy resolves once observation has earned the resolution, and why fixed points are exactly the configurations where this descent finally stops. It does not explain the fountain. A field that only ever flows downhill in energy, monotonically, has no way to sustain the fountain's continuous, patterned motion, because pure descent always terminates, eventually, at a minimum, and a fountain that terminated would not be a fountain anymore, only a puddle. Something in this framework has to account for motion that persists indefinitely without draining away, and gradient descent, taken alone, is structurally incapable of producing it.

\section{From Energy Descent to Phase Space}

The resolution is not to abandon the energy functional Chapter 19 developed, but to use it differently. Rather than always flowing in the direction that decreases H fastest, a field can instead flow along a path that holds H constant, trading one component for another rather than shedding energy at all. This requires treating the field's variables in pairs, a coherence-like quantity and a flow-like quantity bound together the way position and momentum are bound together in ordinary mechanics: Φ plays the role of position, the configuration a point currently occupies, and v plays the role of momentum, the rate and direction at which that configuration is currently changing. A system governed this way does not sit still once it reaches a level set of constant energy. It moves along that level set indefinitely, Φ and v continuously exchanging value with one another, coherence briefly dipping as flow rises to compensate and rising again as flow subsides, forever, in principle, without ever losing energy to the exchange. This is what a fountain is doing, described in the vocabulary this book has been building since Chapter 7.

\section{The Semantic Hamiltonian}

The precise statement of this exchange is a pair of coupled equations, using the same energy functional H from Chapter 19 but now generating motion rather than merely measuring it:

\[ \partial\Phi/\partial t = \delta H/\delta v, \qquad \partial v/\partial t = -\delta H/\delta\Phi. \]

These are Hamilton's equations, applied to the semantic field rather than to a physical particle, and their defining property is conservation: a trajectory that satisfies them keeps H exactly constant along its path, by construction, since each variable's rate of change is defined in terms of the other's contribution to H rather than in terms of H's own gradient with respect to itself. Where Chapter 19's dynamics always pointed downhill, these equations point sideways, along the contour of whatever energy level the field currently occupies, and a field governed purely by this pair of equations oscillates rather than settles. The fountain's water is behaving exactly this way: its appearance and flow trade off against each other continuously, tracing out a stable, repeating pattern rather than draining toward the fountain's minimum-energy configuration, which would simply be still water at rest.

\section{Reconciling Descent and Oscillation}

It would be a mistake to treat this chapter as replacing Chapter 19 rather than completing it, and it is worth being precise about how the two coexist rather than compete. Chapter 17 already established that the field's components operate on separate clocks, fast components like appearance and flow cycling frequently, slow components like coherence and topological structure changing rarely. The Hamiltonian dynamics this chapter introduces govern the fast clock: within a single slow-timescale interval, z and v can trade energy back and forth through many complete cycles, tracing the oscillatory behavior a fountain exhibits from one moment to the next. Chapter 19's gradient descent governs the slow clock: across many such fast-timescale cycles, whatever dissipation the system does experience, however small, accumulates and gradually pulls the slow-timescale coherence Φ toward one of its own minima, exactly the process Chapter 17 described as fast-timescale activity eventually forcing slow-timescale change. A fountain, on this reading, is not purely Hamiltonian and not purely dissipative. It is Hamiltonian at the fast timescale, cycling indefinitely without apparent loss, riding on top of a slow, dissipative drift that governs whether the fountain as a whole is, on the scale of years rather than moments, gradually settling, eroding, or otherwise moving toward a different overall configuration. The two chapters describe the same field at two different rates, not two competing accounts of the same rate.

\section{What the Oscillation Conserves}

It is worth naming, briefly, what exactly stays fixed while a fast-timescale trajectory cycles, because this book has been building toward an answer to this question since Chapter 16 first proposed that identity is a fixed point rather than a label. A trajectory obeying Hamilton's equations conserves more than H alone; it conserves the underlying structure, called symplectic in the branch of geometry this material draws on, that relates Φ and v to one another at every point along the path. This structure is what makes the trajectory recognizably one continuous thing rather than a sequence of unrelated states that happen to share an energy value. The fountain's water, at any given instant, is not the water that occupied the fountain a minute ago, and its specific appearance is different from moment to moment in exactly the way Chapter 9 already described fast-timescale appearance as free to be. What persists is not the water but the trajectory the water is tracing, the conserved structure binding each instant's Φ and v to the next. This refines Chapter 16's claim rather than replacing it: identity, for a genuinely dynamic fixed point like a fountain, is not merely a fixed point of the write cycle in the sense Chapter 16 first stated. It is a closed, symplectically conserved orbit, a fixed pattern of exchange between coherence and flow, and Chapter 16's fixed-point equation turns out to describe the special, degenerate case of this chapter's oscillation where the orbit has shrunk to a single point.

\section{Looking Ahead}

Everything this book's dynamics have described since Chapter 12, writing, projection, commit, fixed points, energy descent, and now conservative oscillation, has been stated pointwise or, at most, across a single trajectory. None of it has yet asked how these local pictures fit together across the field's whole domain at once: whether a coherence structure that looks admissible in one region remains admissible when checked against its neighbors, or whether the affordance-realizability constraint this book has postponed since Chapter 10 can finally be made precise. That is a question about global consistency rather than local dynamics, and it is where Chapter 21 turns next, delivering the sheaf-theoretic bridge this book promised the cue-coherence construction back when affordance was first given its own chapter.
